\documentclass{article}[12pt, a4paper]
\usepackage{amsmath}
\usepackage{fancyhdr}
% \usepackage{lipsum} % For dummy text
\usepackage{xcolor}
% Page style setup
\pagestyle{fancy}
\fancyhf{} % Clear headers and footers
\renewcommand{\headrulewidth}{0pt} % <-- THIS REMOVES THE ANNOYING TOP LINE
\fancyfoot[R]{\textbf{\color{black}em-set-01}}



\begin{document}

% \centering \section*{Quantum Mechanics}

\centering \textbf{(Basics Quantum Mechancis)}
%% Questions
\begin{enumerate}
\item Write down the condition for Orthonormality!
\item Write down the operators for this quantities: $x, p, \hat{H}, E, T, V, L$ \\
($L$ is angular momentum here)
\item Write downt the formulas for: Probability $P(x, t)$, Probability Current Deisnty $J(x)$
Expectation value $\langle f(x) \rangle$ of any function $f(x)$ over a range $(a, b)$
\item Write down the expression for the expectation values of the following quantities
\item Write down uncertinity princple related to $x, p$ and $E, t$
\end{enumerate}

% Basic Quantum Mechanics Problems
\centering \textbf{(Problems)}
\begin{enumerate}
%% Q1
\item A gaussian distribution is given $\rho(x) = A e^{-\lambda (x - a)^2}$ where $A, a$ and $\lambda$ are positive and real constatnts, Find $A$, and then calculate $\langle x \rangle, \langle x^2 \rangle $ and $\sigma$, and sketch the graph of $\rho (x)$

%% Q2
\item At a time $t = 0$, wave function of a particle
$$
\psi (x, 0) =
\begin{cases}
	A \frac{x}{a} & \text{if } 0 \le x \le a, \\
	A \frac{(b - a)}{b - a} & \text{if } a \ge x \ge b, \\
	0 & \text{otherwise}

\end{cases}$$
\begin{enumerate}
	\item Normalize $\psi$
	\item Sketch $\psi(x, 0)$, as function of $x$
\end{enumerate}

%% Q3
\item Consider the wave function $$\Psi(x,t) = Ae^{-\lambda |x|}e^{-i\omega t}$$
where $A, a$ and $b$ are constants
\begin{enumerate}
	\item Normalize $\Psi$
	\item Determine the expetcation values of $x, x^2$
\end{enumerate}

%% Q4
\item The neddle on a broken speedometer is free to swing, and bounces perfectly off the pins at
either end, so that if you give it a flick it is equally likely to come to rest at any angle between $0$ and $\pi$.
\begin{enumerate}
	\item What is the probability density $\rho(\theta)$ ? \textbf{Hint:} $\rho(\theta)d\theta$ is the
	is the probability that the needle will come to rest between $\theta$ and $\theta+d\theta$.
	Graph $\rho(\theta)$ as a function of $\theat$ between $(-\pi/2, 3\pi/2)$ (Of course, part of this interveal excluded
	so $\rho$ is zero there.) Make sure that the total probability is $1$
	\item Compute $\langle \theta \rangle, \langle x^2 \rangle$ and $\sigma$ for this distribution
	\item Compute $\langle sin(\theta) \rangle, \langle cos(\theta) \rangle, \langle cos^2(\theta) \rangle$
\end{enumerate}
\end{enumerate}





\textbf{(One Dimentional Potentials)}
\begin{enumerate}

%% Stationary States
\item Start with General schrodinger equation,
$$\frac{-\hbar^2}{2m}\frac{\partial^2}{\partial x^2} \Psi(x, t) + V(x) \Psi(x, t)= -i \hbar \frac{\partial}{\partial t} \Psi(x, t)$$ \\
Show that for Stationary state solution (fixed $E$ state), general linear solution of the equation is
$$\Psi(x, t) = \sum_n C_n \psi_n(x)exp({\frac{-i E_n t}{\hbar}})$$
\begin{enumerate}
    \item Write down the \textbf{Fourier's trick} to obtain $C_n$ for $nth$ energy state.
    \item Show that for Stationary state, probability density $\mathbf{P}$ independent of time $t$.
	\item Show that for Stationary States flactuation of energy ($E$) is $\sigma_{\hat{H}} = 0$
\end{enumerate}

% Infinify potential Wall
\item Solve \textbf{TISE} for the potential given
$$
V(x) =
\begin{cases}
0 & \text{if } 0 \le x \le a \\
\infty & \text{otherwise}
\end{cases}$$
\begin{enumerate}
 \item Show that total normalized solution is $\psi(x) = \displaystyle\mathbf{\sum_n} C_n \sqrt{\frac{2}{a}} \sin \displaystyle(\frac{\displaystyle n\pi x}{\displaystyle a}\displaystyle)$
 \item Find the expression of energy of $nth$ state $E_n$.
 \item Find $C_n$ using \textbf{Fourier's trick}, wirte down the expression.
\end{enumerate}

% Potential Barrier
\item Solve \textbf{TISE} for the potential given (Potential Step)
$$
V(x) =
\begin{cases}
0 & \text{if }  x \le 0 \\
V_o & \text{if } x \ge 0 \\
\end{cases}$$
\begin{enumerate}
    \item Sketch the potential $V(x)$ for case $(E > V_o)$ and $(V_o > E)$.
	\item For $(E > V_o)$ case solve the \textbf{TISE} and write down the valid solutions and boundary conditions.
	\item For $(E > V_o)$ show that, $T + R = 1$ (where $T$ is the transmission cofficient and $R$ is the reflection cofficient).
	\item Show that when $V_o \gg E$, then $R = 1, T = 0$.
	\item Show that for $(V_o > E)$, transmission cofficient $J_{trans} = 0$, and hence show that
	reflection cofficient $R = 0$
\end{enumerate}

% Potential Barrier
\item For the given one dimentional potential
$$
V(x) =
\begin{cases}
0 & \text{if }  x \le 0 \\
V_o & \text{if } 0 < x < a \\
0 & \text{if }  x \ge a \\
\end{cases}$$
\begin{enumerate}
   \item Sketch the potential $V(x)$ for case $(E > V_o)$ and $(V_o > E)$ and write down the \textbf{boundary conditions}.
   \item
\end{enumerate}
\end{enumerate}

% One Dimentional Potential Problems
\centering \textbf{(Problems)}
\begin{enumerate}

\item Solve the time independent Schroedinger Equation with appropiate boundary conditions
for infinity square well potential, centered at the origin [$V(x) = 0$ for $x \in (-a/2, a/2)$ and $V = \infty$ otherwise]
Calculate $\langle x \rangle, \langle x^2 \rangle, \langle p \rangle, \langle p^2 \rangle, \sigma_x$ and $\sigma_p$,
for it's $n^{th}$ stationary state. Check the \textbf{uncertinity principle} is satisfied. Which state
comes close to the uncertinity limit?

\item Consider the double delta function potential $$V(x) = -\alpha[\delta(x + a) + \delta(x - a)],$$
where $\alpha, a$ constant.
\begin{enumerate}
	\item Sketch the potential
	\item How many bound state does it possess? Find the allowed energies for, $\alpha = \hbar^2/ma$ and for
	$\alpha = \hbar^2/4ma$, and sketch the wave functions
	\item Find the transmission cofficient for this potential.
 \end{enumerate}

\end{enumerate}


\centering \textbf{(3 Dimentional Potential Problems)}
\begin{enumerate}
\item Write down



\end{enumerate}


\end{enumerate}


\end{document}
